% ===================================================================
%  TABELO N-RO 8 — La rikolto. — La ĉasado, la vinberrikolto, la
%  fiŝkaptado.
%  Traduction 2026 d'apres texto/io, controlee sur texto/fr et
%  texto/es. Consigne : voir texto/eo/10-tabelo-01.tex.
%
%  L'appel de note est dans le TITRE de la scene — « La rikolto (*) »
%  — et non dans un alinea : c'est ainsi que l'ido le compose.
%
%  LE RENVOI (13) DE LA FAUX MANQUE A L'IDO et le francais le donne.
%  La colonne esperantiste le rend, comme les autres traductions :
%  c'est le deuxieme des trois ecarts declares dans renvoji.py.
% ===================================================================

%%K t08-noto-1 noto
\VUnotes{143.2mm}{%
\textit{(*) Ĉar tiu vorto estas ne preciza: ĉu estas linrikolto,
vinberrikolto, pomrikolto? Eble estus bone uzi «} \VUgras{mesono}
\textit{» proponitan en Mondo XIV, p. 242. Sed ĝis nun tiu nova radiko ne estas
adoptita de la Akademio kaj atendas diskutojn laŭ la kutima regularo idista.}%
}

%%K t08-apar-1 apar
\VUsaut{5.0mm}
\VUtitre{15.0pt}{60}{TABELO N\textsuperscript{o} 8}
\VUsaut{3.0mm}
\VUcentre{13.2pt}{90}{\textit{Unua sceno.}}
\VUsaut{3.0mm}
\VUpk{10.2pt}{40}{La rikolto (*).}
\VUsaut{4.0mm}
\VUpk{10.2pt}{40}{La aspektoj de la kamparo.}
\VUsaut{3.0mm}

%%K t08-c1-01-1 p
1. --- Kun la \VUgras{somero}, unu plian fojon ŝanĝiĝis la aspekto de la
kamparo. La semo konfidita al la sulko ĝermis; verda plantaĵo eliris el la tero
kaj spikiĝis. La grajno estis formata, poste ĝi maturiĝis kaj eĉ de nun oni
bezonas nur rikolti.

%%K t08-c1-02-1 p
2. --- La \VUgras{kamparaj floroj} estas malfermitaj. \VUgras{Cejanoj}
\textsuperscript{(1)}, \VUgras{papavetoj} \textsuperscript{(2)} kaj
\VUgras{digitaloj} \textsuperscript{(3)} miksas al la unuforma nuanco de la
tritikkampoj la gajan kontraston de siaj freŝaj \VUgras{koroloj}.

%%K t08-c1-03-1 p
3. --- La fruktarboj sin garnis per folioj; la frukto maturiĝis. La malgranda
\VUgras{ĉerizarbo} \textsuperscript{(4)} abundas je belaj \VUgras{ĉerizoj}
\textsuperscript{(5)} kiujn la knaboj plezurege kolektas kaj manĝas.

%%K t08-c1-04-1 p
4. --- De nun la brutaro povas pasigi la tutan tagon en la libera aero.

%%K t08-c1-04-2 p
La \VUgras{ŝafidoj} \textsuperscript{(6)} grandiĝis kaj fariĝis belaj
\VUgras{ŝafinoj} \textsuperscript{(7)} kaj belaj \VUgras{virŝafoj}
\textsuperscript{(8)} kun densa lanfelo. Ili trankvile paŝtas la \VUgras{herbon}
de la \VUgras{paŝtejo} \textsuperscript{(9)} buntigita per \VUgras{lekantetoj}.

%%K t08-c1-04-3 p
Baldaŭ oni ektondos tiujn bestojn por havi ilian \VUgras{lanon} el kiu oni
fabrikas la drapon de niaj vestoj.

%%K t08-tit-2 sub
\VUsaut{5.0mm}
\VUpk{10.2pt}{40}{La paŝtisto.}
\VUsaut{3.4mm}

%%K t08-c1-05-1 p
5. --- La \VUgras{paŝtisto} \textsuperscript{(10)} ŝajnas ne tre atenti ilin. Li
zorgas nur pri la ŝtormo kiu pasas ĉe la horizonto, kaj la \VUgras{ŝafgardisto}
\textsuperscript{(13)}, kiu proksimigas \VUgras{kruĉon} \textsuperscript{(14)}
al siaj lipoj ŝajnas eĉ pli senzorga.

%%K t08-c1-06-1 p
6. --- La varmego estas premanta; ĉar la \VUgras{hundo} \textsuperscript{(15)}
venas, ĝi ankaŭ, sensoifigi sin; ĝi lekas la freŝan akvon de la
\VUgras{rivereto} \textsuperscript{(16)} kaj timigas kompatindan
\VUgras{raneton} \textsuperscript{(17)} kiu kaŝiĝis en la bordoherboj. Ĉi tiu
saltas en la klaran akvon kie floras la \VUgras{nimfeoj}
\textsuperscript{(18)}.

%%K t08-c1-07-1 p
7. --- \VUgras{Motacilo} \textsuperscript{(19)} sin banas apud la
\VUgras{fonteto} \textsuperscript{(20)} kiu ŝprucas inter du rokaĵoj kaj sin
kaŝas sub la \VUgras{dornaj arbustoj} \textsuperscript{(21)} kaj la
\VUgras{kratagaj arbetoj} \textsuperscript{(22)}.

%%K t08-tit-3 sub
\VUsaut{5.0mm}
\VUpk{10.2pt}{40}{La rikolto.}
\VUsaut{3.4mm}

%%K t08-c1-08-1 p
8. --- Por la kamparaj knaboj, la tritikrikolto estas tre amuza epoko. Ili ĝoje
\VUgras{transkapiĝas} \textsuperscript{(24)} sur la \VUgras{pajlamasoj}
\textsuperscript{(23)}, ili helpas la \VUgras{rikoltistojn}
\textsuperscript{(26)} kaj dum ĉi tiuj falĉas la \VUgras{tritikkampon}
\textsuperscript{(27)} per siaj \VUgras{falĉiloj} \textsuperscript{(25)} bone
akrigitaj per la \VUgras{akriga ŝtono} \textsuperscript{(30)}, ili kolektas la
tritikon kaj ligas ĝin kiel \VUgras{garbojn} \textsuperscript{(31)} aŭ kiel
\VUgras{fasketojn} \textsuperscript{(32)}.

%%K t08-c1-09-1 p
9. --- Kelkfoje oni permesas al la plej grandaj inter ili tranĉi per
\VUgras{rikoltilo} \textsuperscript{(12)} la \VUgras{orajn tigojn}
\textsuperscript{(28)} kiuj portas la pezajn \VUgras{spikojn}
\textsuperscript{(29)} grajnoplenajn; kaj la malgrandaj, gvatante la
\VUgras{brutaron} \textsuperscript{(33)} kiu paŝtas en la \VUgras{paŝtejo}
\textsuperscript{(34)} sub la ombro de la granda \VUgras{tilio}
\textsuperscript{(35)} kaptas per sia \VUgras{reto} \textsuperscript{(36)} la
belajn \VUgras{papiliojn}.

%%K t08-c1-09-2 p
Eĉ kelkaj grimpas sur la \VUgras{ĉareton} \textsuperscript{(37)} por ordigi la
garbojn kiujn oni transdonas al ili per \VUgras{forko}
\textsuperscript{(38)}.

%%K t08-c1-10-1 p
10. --- Sur la tuta \VUgras{ebenaĵo} \textsuperscript{(39)} kie ekflugas la
\VUgras{alaŭdoj} \textsuperscript{(41)} regas la vivo. Ĉiuj estas okupataj pri
la rikolto. Sed, dum la malriĉuloj daŭrigas uzi la malnovajn instrumentojn,
\VUgras{falĉilojn, rikoltilojn} kaj \VUgras{draŝilojn}
\textsuperscript{(40)}, la riĉaj bienposedantoj falĉigas sian tritikon aŭ sian
\VUgras{sekalon} \textsuperscript{(43)} per \VUgras{falĉmaŝino}
\textsuperscript{(42)}. Poste, ne bezonante porti la garbojn en la garbejon,
oni faras grandajn \VUgras{garbamasojn} \textsuperscript{(44)}.

%%K t08-c1-11-1 p
11. --- Tiam oni alkondukas \VUgras{draŝmaŝinon} \textsuperscript{(47)} kiun
\VUgras{lokomobilo} \textsuperscript{(45)} movas per \VUgras{transmisia rimeno}
\textsuperscript{(46)} kaj tiel la grajno separiĝas de la \VUgras{pajlo}, sen
forlasi la kampon kie ĝi estis semata.

%%K t08-tit-4 sub
\VUsaut{5.0mm}
\VUpk{10.2pt}{40}{La ŝtormo.}
\VUsaut{3.4mm}

%%K t08-c1-12-1 p
12. --- Ofte \VUgras{ŝtormegoj} \textsuperscript{(53)} okazas, dum la tempo de
la rikolto. Unue la grumblado de la \VUgras{tondro} estas aŭdata de
malproksime; \VUgras{ventego el okcidento} \textsuperscript{(54)} ekblovas kaj
anhelante kurbigas la plej dikajn arbojn de la \VUgras{arbaro}
\textsuperscript{(49)}. Nigraj \VUgras{nubegoj} \textsuperscript{(56)} atakas la
ĉielon; ili estas trairataj de zigzagoj blindigaj de la \VUgras{fulmo}
\textsuperscript{(55)}. La \VUgras{pluvo} kaj kelkfoje la \VUgras{hajlo}
ekfalas, platigante la rikoltaĵojn pretajn por esti falĉataj.

%%K t08-c1-13-1 p
13. --- Ofte la fulmo incendias konstruaĵon. Tiel lastjare la tegmento de la
\VUgras{ventmuelilo} \textsuperscript{(50)} estis cindrigita kaj du el ĝiaj
\VUgras{flugilegoj} \textsuperscript{(51)} frakasitaj.

%%K t08-c1-14-1 p
14. --- Post kelkaj horoj, la kvieto renaskiĝas. Brila \VUgras{ĉielarko}
\textsuperscript{(52)} aperas ĉe la horizonto kaj la naturo reprenas sian
vivon interrompitan dum kelkaj momentoj. La kamparanoj, kiuj rifuĝis en la
\VUgras{kabanojn} \textsuperscript{(48)} relaboras kaj kuraĝe penas ripari la
damaĝojn kiujn ili ĵus suferis.

%%K t08-tit-5 sub
\VUsaut{6.0mm}
\VUcentre{13.2pt}{90}{\textit{Dua sceno.}}
\VUsaut{3.6mm}

%%K t08-tit-6 sub
\VUpk{10.2pt}{40}{La ĉasado. --- La vinberrikolto.}
\VUsaut{1.4mm}
\VUpk{10.2pt}{40}{La fiŝkaptado.}
\VUsaut{4.0mm}

%%K t08-c2-01-1 p
1. --- Ĉirkaŭ la fino de \VUgras{aŭgusto} aŭ la mezo de \VUgras{septembro}, laŭ
la regionoj, okazas la malfermo de la ĉasado. Apenaŭ la unuaj sunradioj igis la
\VUgras{lacertojn} \textsuperscript{(2)} eliri el sia truo, kiam la
\VUgras{ĉasisto} \textsuperscript{(5)} ekvojiras kun siaj \VUgras{hundoj}
\textsuperscript{(4)}.

%%K t08-c2-02-1 p
2. --- Li metis sian solidan \VUgras{ĉaskostumon} \textsuperscript{(9)} el dika
drapo, piedvestis sin per ledaj \VUgras{gamaŝoj}, per kiuj li sukcesos marŝi en
la humida herbo kie sin kaŝas la \VUgras{bufoj} \textsuperscript{(1)}; li prenis
sian \VUgras{fusilon} \textsuperscript{(6)} kaj sian \VUgras{ĉassakon}
\textsuperscript{(10)}, kaj jen ke li foriras.

%%K t08-c2-03-1 p
3. --- La fervoro de la persekuto alkondukas lin meze de vitkampo kie la
vinberrikolto estas tre vigla. La knaboj de la vitkultivisto, kiuj ordinare
gvatas la kaprinojn sur la vojo, hodiaŭ lasas ilin paŝti hazarde, kaj post
ligi al paliso maljunan \VUgras{virkapron} \textsuperscript{(3)} kiu ne
mankus profiti sian liberecon por eskapi, ili, siavice, atribuas al si
partoprenon en la plezuro. Unu ekprenas vinberaron en \VUgras{rikoltokorbo}
\textsuperscript{(12)} kaj amuzas sin fluigi la \VUgras{sukon}
\textsuperscript{(14)} en pokaleton \textsuperscript{(13)} por fabriki moston
(vinon ne fermentintan). Alia kombas sin per \VUgras{foliokrono}
\textsuperscript{(15)} kiun li ĵus plektis. Apud ili, senhontaj
\VUgras{paseroj} \textsuperscript{(16)} venas eĉ antaŭ la premilon por rabi la
vinberojn.

%%K t08-c2-04-1 p
4. --- Dum la knaboj sin amuzas, la viroj laboras. La
\VUgras{vinberrikoltistoj} \textsuperscript{(18)} alportas la vinberojn en la
kelon per \VUgras{dorskorboj} \textsuperscript{(17)}; \VUgras{barelfaristo}
\textsuperscript{(19)} riparas la \VUgras{barelojn} \textsuperscript{(20)},
kontrolas ĉu la \VUgras{duvoj} \textsuperscript{(22)} kaj la \VUgras{fundoj}
\textsuperscript{(21)} estas bonstataj, kaj anstataŭigas la \VUgras{ringegojn}
\textsuperscript{(23)} rompitajn. Li ankaŭ laboris la \VUgras{kuvegojn}
\textsuperscript{(25)} en kiujn oni verŝas la \VUgras{vinberojn}
\textsuperscript{(26)} por igi ilin fermenti.

%%K t08-c2-05-1 p
5. --- Kiam la fermentado finiĝis, oni ĉerpas la vinon kaj oni metas la
restaĵon sur la \VUgras{premilon} \textsuperscript{(28)} kaj progresive
premante per la \VUgras{ŝraŭbego} \textsuperscript{(29)} oni elpremas la sukon
kiu fluas en \VUgras{kuveton} \textsuperscript{(32)} kaj ankoraŭ oni akiras
\VUgras{vinon} \textsuperscript{(31)}.

%%K t08-tit-8 sub
\VUsaut{6.0mm}
\VUpk{10.2pt}{40}{Biero kaj cidro.}
\VUsaut{3.4mm}

%%K t08-c2-06-1 p
6. --- Ĉe la muro de la \VUgras{kelo} \textsuperscript{(27)} grimpas branĉo de
\VUgras{lupolo} \textsuperscript{(30)}. Tie ĝi estas nur ornama planto, sed en
kelkaj landoj, kie la vito estas tre malofta, la lupolo estas kultivata por
fabriki la \VUgras{bieron}.

%%K t08-c2-07-1 p
7. --- La malgranda \VUgras{pomarbo} \textsuperscript{(33)} estas ŝarĝita per
pomoj kiuj produktos kiam ili estos maturaj, bonegan \VUgras{cidron}. En sia
\VUgras{kaĝo} \textsuperscript{(34)} la \VUgras{kanario} \textsuperscript{(35)}
kaj la \VUgras{kardelo} \textsuperscript{(36)} rigardas per enviaj okuloj la
paserojn kiuj libere petolas.

%%K t08-c2-08-1 p
8. --- Ĝis la fino de la vido, sur la deklivo de la montetoj, estas viciĝintaj
\VUgras{vitstangoj} \textsuperscript{(40)} kiuj subtenas la belegajn
\VUgras{(vit)trunkojn} \textsuperscript{(39)}, garnitajn de \VUgras{beraroj}
pezaj.

%%K t08-c2-08-2 p
Dum la tuta jaro, la \VUgras{vitkultivisto} \textsuperscript{(42)} laboris por
zorgi sian \VUgras{vitkampon} \textsuperscript{(38)} kaj nun li estas
rekompencata pro sia peno per bela rikolto. La \VUgras{vinberrikoltistinoj}
\textsuperscript{(41)} plenigas la kuvojn metitajn sur la ĉareton tiratan de
\VUgras{muloj} \textsuperscript{(43)}, kaj ĉiuj estas ĝojaj.

%%K t08-tit-9 sub
\VUsaut{5.0mm}
\VUpk{10.2pt}{40}{La fiŝkaptado.}
\VUsaut{3.4mm}

%%K t08-c2-09-1 p
9. --- Same estas pri la \VUgras{fadenfiŝistoj} \textsuperscript{(44)} kiuj de
la mateno, venis instali sin ĉe la bordo de la akvo kun siaj
\VUgras{fiŝkaptobastonoj} \textsuperscript{(45)}.

%%K t08-c2-09-2 p
La \VUgras{fiŝoj} \textsuperscript{(47)} kaptas la \VUgras{hokon} tuj kiam ili
plonĝas sian \VUgras{fadenon} \textsuperscript{(46)} en la akvon, kaj apenaŭ
ili havis tempon por renovigi la \VUgras{logaĵon}, kiam nova viktimo baraktas
ĉe la ekstremo de la fadeno.

%%K t08-c2-09-3 p
Tamen la \VUgras{retfiŝisto} \textsuperscript{(48)} kiu, el sia \VUgras{barko}
\textsuperscript{(49)}, ĵetas la \VUgras{(ĵet)reton} en la mezon de la
\VUgras{rivero} \textsuperscript{(50)} kaptas pli multajn fiŝojn.

%%K t08-c2-10-1 p
10. --- La aŭtuno estas la sezono de la bonaj promenadoj: piede,
\VUgras{rajde} \textsuperscript{(52)}, \VUgras{bicikle} \textsuperscript{(53)},
\VUgras{aŭtomobile} \textsuperscript{(54)}. La \VUgras{vojoj}
\textsuperscript{(51)} estas puraj kaj oni profitas la lastajn belajn tagojn,
ĉar jen ke jam la \VUgras{hirundoj} \textsuperscript{(56)} kunigas sin sur la
tegmentoj por ekflugi per ĉiuj flugiloj al pli varmaj landoj. La foliaro de la
maljuna \VUgras{juglandarbo} \textsuperscript{(57)} flaviĝas, la
\VUgras{juglandoj} falas kaj la altaj \VUgras{arboj} grupe dispoziciitaj kiel
\VUgras{bosketo} \textsuperscript{(58)} sur la \VUgras{altaĵo}
\textsuperscript{(59)} baldaŭ senfoliiĝos.

%%K t08-c2-11-1 p
11. --- La \VUgras{suno} \textsuperscript{(60)} sin levas pli malfrue, la tagoj
malkreskas, malvarmeto sufiĉe pika komencas esti sentata, ĉe la \VUgras{mateno},
kaj jen ke jam flugoj de \VUgras{skolopoj} \textsuperscript{(61)} forlasas niajn
klimatojn negastemajn.
\VUsaut{6.0mm}
\VUfilet{22mm}
